\documentclass[11pt]{article}
\usepackage[margin=0.9in]{geometry}
\usepackage{amsmath,amssymb,amsthm,mathtools}
\usepackage[T1]{fontenc}
\usepackage[utf8]{inputenc}
\usepackage{enumitem}
\usepackage{booktabs,longtable,array,seqsplit}
\usepackage{hyperref}
\usepackage{xcolor}
\usepackage{microtype}
\hypersetup{colorlinks=true,linkcolor=blue!45!black,citecolor=blue!45!black,urlcolor=blue!45!black}
\setlength{\emergencystretch}{6em}

\newtheorem{theorem}{Theorem}[section]
\newtheorem{proposition}[theorem]{Proposition}
\newtheorem{lemma}[theorem]{Lemma}
\newtheorem{corollary}[theorem]{Corollary}
\newtheorem{claim}[theorem]{Claim}
\theoremstyle{definition}
\newtheorem{definition}[theorem]{Definition}
\newtheorem{assumption}[theorem]{Assumption}
\theoremstyle{remark}
\newtheorem{remark}[theorem]{Remark}

\newcommand{\R}{\mathbb{R}}
\newcommand{\C}{\mathbb{C}}
\newcommand{\N}{\mathbb{N}}
\newcommand{\Hsp}{\mathcal{H}}
\newcommand{\Fsp}{\mathcal{F}}
\newcommand{\Id}{\mathrm{Id}}
\newcommand{\tr}{\operatorname{tr}}
\newcommand{\diag}{\operatorname{diag}}
\newcommand{\rank}{\operatorname{rank}}
\newcommand{\spec}{\operatorname{spec}}
\newcommand{\spanop}{\operatorname{span}}
\newcommand{\dist}{\operatorname{dist}}
\newcommand{\range}{\operatorname{ran}}
\newcommand{\Ker}{\operatorname{ker}}
\newcommand{\Pinch}{\mathcal{C}}
\newcommand{\Off}{\widetilde{\mathcal{C}}}
\newcommand{\Pperp}{P^{\perp}}
\newcommand{\Qperp}{Q^{\perp}}
\newcommand{\op}{\mathrm{op}}
\newcommand{\F}{\mathrm{F}}
\newcommand{\HS}{\mathrm{HS}}
\newcommand{\ip}[2]{\left\langle #1,#2\right\rangle}
\newcommand{\norm}[1]{\left\lVert #1\right\rVert}
\newcommand{\abs}[1]{\left\lvert #1\right\rvert}
\newcommand{\code}[1]{\nolinkurl{#1}}

\newenvironment{argument}{\paragraph{Argument.}\begin{enumerate}[leftmargin=2em]}{\end{enumerate}}
\newenvironment{proofoutline}{\paragraph{Proof architecture.}\begin{enumerate}[leftmargin=2em]}{\end{enumerate}}
\newcommand{\sourceanchor}[1]{\par\smallskip\noindent\textbf{Source anchor.} #1\par\smallskip}
\newcommand{\sourcecomment}[1]{\begin{remark}#1\end{remark}}
\newenvironment{sourceguide}{\begin{description}[style=nextline,leftmargin=3.3cm,labelwidth=3.1cm]}{\end{description}}
\newenvironment{formalizationmap}{\begin{longtable}{@{}p{0.20\textwidth}p{0.31\textwidth}p{0.40\textwidth}@{}}\toprule Source item & Modern mathematical content & Repository status \\\midrule\endfirsthead\toprule Source item & Modern mathematical content & Repository status \\\midrule\endhead}{\bottomrule\end{longtable}}

\newcommand{\M}{\mathcal{M}}
\newcommand{\Z}{\mathbb{Z}}
\newcommand{\supp}{\operatorname{supp}}
\newcommand{\sgn}{\operatorname{sgn}}
\newcommand{\TV}{\operatorname{TV}}
\newcommand{\cF}{\mathcal{F}}
\newcommand{\Lip}{\operatorname{Lip}}
\newcommand{\esssup}{\operatorname*{ess\,sup}}
\renewenvironment{formalizationmap}{%
  \begin{longtable}{@{}p{0.20\textwidth}p{0.31\textwidth}p{0.40\textwidth}@{}}%
  \toprule Source item & Modern mathematical content & Repository status \\\midrule\endhead%
}{\bottomrule\end{longtable}}
\title{Source-order modernization of Haagerup--Zsid\'o (1994): sharp Hermitian resolvents and minimal Fourier extrapolation}
\author{}
\date{Distilled reconstruction from the supplied published article}
\begin{document}
\maketitle

\section*{Source map and scope}
\begin{sourceguide}
\item[Primary source] Uffe Haagerup and L\'aszl\'o Zsid\'o, \emph{Resolvent Estimate for Hermitian Operators and a Related Minimal Extrapolation Problem}, Acta Scientiarum Mathematicarum (Szeged) 59 (1994), 503--524.
\item[Local evidence] Supplied PDF \code{595055268644.pdf}.\\ SHA--256 \texttt{\seqsplit{77fe9c625a913bcd46b9b0038aa3321039c777fce0d35e45e738b0faabd9642c}}.
\item[Redistribution] The source PDF is not vendored. This document is a new mathematical reconstruction, correction ledger, and formalization map.
\item[Coverage] Introduction; Theorems 1.1, 1.2, 1.5; Lemma 1.3; Corollaries 1.4 and 1.6; Theorem 2.1; Corollaries 2.2 and 2.4; Theorem 2.3; Propositions 3.1 and 3.2; Lemmas 3.3--3.5; Theorem 3.6; all displayed identities used in their proofs; the sharpness construction; and the source's two typographical slips relevant to formalization.
\item[Modernization] Fourier-transform conventions, dual-pairing arguments, operator-valued integration, support statements, and scaling are made explicit. ``Hermitian'' is kept in the Banach-space generator sense used by the paper and is not conflated with Hilbert-space self-adjointness.
\item[Formalization target] The scalar analytic root in
\texttt{\seqsplit{DavisKahan/Alternative/FiniteDimensional/Sylvester/ContinuousLinearMapBridge.lean}}, especially
\texttt{\seqsplit{sylvester\_barycentricOrbitRepresentation\_of\_spectralDistance}}. The special case $\alpha=0$ of Theorem 3.6 gives the sharp reciprocal Fourier multiplier of total $L^1$ mass $\pi/2$ needed by that theorem.
\end{sourceguide}

This note follows the source in proof order. It does not merely quote the final resolvent bound: the full abstract minimal-extrapolation theorem, the sign certificate, the explicit extremizer, the norm computation, and the passage to operator resolvents are reconstructed because each is a possible formalization seam.

\section{Conventions and modern terminology}

\subsection{Fourier convention}

For $f\in L^1(\R)$ the paper uses the inverse-Fourier-sign convention
\[
  \widehat f(x):=\int_{\R} f(t)e^{ixt}\,dt.
  \tag{HZ-F}\label{HZ-F}
\]
For a bounded complex regular Borel measure $\mu\in M(G)$ on a locally compact abelian group $G$,
\[
  \widehat\mu(\gamma):=\int_G \gamma(t)\,d\mu(t),
  \qquad \gamma\in\widehat G.
  \tag{HZ-M}\label{HZ-M}
\]
An $L^1$ function is identified with its absolutely continuous measure $f(t)\,dt$.

\subsection{Hermitian operators on Banach spaces}

Let $X$ be a complex Banach space and let $T$ be a densely defined closed operator. The paper calls $T$ \emph{Hermitian} when either, hence both, of the following equivalent statements hold:
\begin{enumerate}[label=(\roman*)]
\item $\sigma(T)\subset\R$ and
\[
  \norm{(\lambda I-T)^{-1}}
  \le \frac{1}{\abs{\Im\lambda}}
  \qquad(\lambda\in\C\setminus\R);
  \tag{HZ-H1}\label{HZ-H1}
\]
\item $iT$ is the infinitesimal generator of a strongly continuous one-parameter group $(U_t)_{t\in\R}$ of linear contractions. We record this as
\[
  iT=\operatorname{gen}(U_t)
  \quad\text{with }\norm{U_t}\le1\text{ for all }t\in\R.
  \tag{HZ-H2}\label{HZ-H2}
\]
\end{enumerate}
Because $U_tU_{-t}=I$ and both $U_t$ and $U_{-t}$ are contractions, each $U_t$ is in fact an isometry.

On a Hilbert space, a self-adjoint operator satisfies the stronger exact identity
\[
  \norm{(\lambda I-T)^{-1}}
  =\frac{1}{\dist(\lambda,\sigma(T))}.
  \tag{HZ-Hilbert}\label{HZ-Hilbert}
\]
The paper determines the best universal replacement for general Banach-space Hermitian operators.

\section{Introduction and the inherited estimates}

\sourceanchor{Introduction and Section 1 through Corollary 1.4.}

Partington had shown that a universal constant exists and that it cannot be smaller than $\pi/2$.

\begin{theorem}[Partington, source Theorem 1.1]
Let $T$ be Hermitian.
\begin{enumerate}[label=(\arabic*)]
\item For every real $\lambda\notin\sigma(T)$,
\[
  \norm{(\lambda I-T)^{-1}}
  \le \frac{\pi}{2}\,\dist(\lambda,\sigma(T))^{-1},
  \tag{HZ-1.1a}\label{HZ-1.1a}
\]
and $\pi/2$ is optimal.
\item For every $\lambda\in\C\setminus\sigma(T)$,
\[
  \norm{(\lambda I-T)^{-1}}
  \le \sqrt{1+\frac{\pi}{2}}\,
      \dist(\lambda,\sigma(T))^{-1}.
  \tag{HZ-1.1b}\label{HZ-1.1b}
\]
\end{enumerate}
\end{theorem}

The conjecture settled by Haagerup--Zsid\'o is that the constant in the second statement can also be $\pi/2$.

\begin{theorem}[Zsid\'o, source Theorem 1.2]
Let $T$ be closed, let $\sigma(T)\subset\R$, and suppose that for constants $m,c\ge0$,
\[
  \norm{(\lambda I-T)^{-1}}
  \le c\abs{\Im\lambda}^{-(m+1)}
  \qquad(\lambda\in\C\setminus\R).
  \tag{HZ-1.2-hyp}\label{HZ-1.2-hyp}
\]
Then
\[
  \norm{(\lambda I-T)^{-1}}
  \le (2\sqrt e)^{m+1}c\,
      \dist(\lambda,\sigma(T))^{-(m+1)}
  \qquad(\lambda\in\C\setminus\sigma(T)).
  \tag{HZ-1.2}\label{HZ-1.2}
\]
\end{theorem}

The analytic input is the following subharmonic estimate.

\begin{lemma}[Source Lemma 1.3]
Let $\Omega\subset\C$ be open and let $u:\Omega\to[-\infty,+\infty)$ be subharmonic. If
\[
  u(\lambda)\le \log\abs{\Im\lambda}^{-1}
  \qquad(\lambda\in\Omega\setminus\R),
  \tag{HZ-1.3-hyp}\label{HZ-1.3-hyp}
\]
then
\[
  u(\lambda)
  \le \log\!\left(2\sqrt e\,\dist(\lambda,\partial\Omega)^{-1}\right)
  \qquad(\lambda\in\Omega).
  \tag{HZ-1.3}\label{HZ-1.3}
\]
\end{lemma}

Partington later reduced $2\sqrt e$ to the optimal constant $2$ for this lemma. For $m=0$ and $c=1$ this gives:

\begin{corollary}[Source Corollary 1.4]
Every Hermitian operator satisfies
\[
  \norm{(\lambda I-T)^{-1}}
  \le 2\,\dist(\lambda,\sigma(T))^{-1}
  \qquad(\lambda\in\C\setminus\sigma(T)).
  \tag{HZ-1.4}\label{HZ-1.4}
\]
\end{corollary}

The source also records that a more elaborate subharmonic argument lowers $2$ to an approximately $1.87$ constant, but not to the conjectured sharp $\pi/2$.

\section{The normalized resolvent theorem and its sharpness example}

\sourceanchor{Section 1, Theorem 1.5 and Corollary 1.6.}

\begin{theorem}[Source Theorem 1.5]
Let $T$ be Hermitian and suppose
\[
  \sigma(T)\subset\R\setminus(-1,1).
\]
Then for every $\alpha\in\R$,
\[
  \norm{(i\alpha I+T)^{-1}}
  \le
  \begin{cases}
    \displaystyle \frac{1}{\alpha}\tanh\!\left(\frac{\pi\alpha}{2}\right),
      & \alpha\ne0,\\[1.1ex]
    \displaystyle \frac{\pi}{2},&\alpha=0.
  \end{cases}
  \tag{HZ-1.5}\label{HZ-1.5}
\]
The expression for $\alpha\ne0$ is positive because $\tanh$ is odd.

Moreover, there is a Hermitian $T$ with
\[
  \sigma(T)=2\Z+1
\]
for which equality holds for every $\alpha\in\R$.
\end{theorem}

\subsection{Upper bound from the extrapolating kernel}

Let $(U_t)$ be the isometric group generated by $iT$. For $\alpha>0$, the Hille--Yosida resolvent formula gives, in the weak-operator/Pettis sense,
\[
  i(i\alpha I+T)^{-1}
  =(\alpha I-iT)^{-1}
  =\int_0^\infty e^{-\alpha t}U_t\,dt.
  \tag{HZ-1.5-res}\label{HZ-1.5-res}
\]
The integrated group representation
\[
  \widehat f\longmapsto \int_{\R}f(t)U_t\,dt
  \tag{HZ-fc}\label{HZ-fc}
\]
has support equal to $\sigma(T)$. Hence, whenever $f,g\in L^1(\R)$ satisfy
\[
  \widehat f(x)=\widehat g(x)
  \qquad\text{for every }x\in\R\setminus(-1,1),
  \tag{HZ-support}\label{HZ-support}
\]
their integrated operators agree.

Theorem 3.6 constructs $f_\alpha\in L^1(\R)$ with
\[
  \widehat f_\alpha(x)=\frac{1}{\alpha-ix}
  \quad(|x|\ge1),
  \qquad
  \norm{f_\alpha}_1
  =\frac1\alpha\tanh\!\left(\frac{\pi\alpha}{2}\right).
  \tag{HZ-ext-alpha}\label{HZ-ext-alpha}
\]
Since $e^{-\alpha t}\mathbf1_{[0,\infty)}(t)$ has Fourier transform $(\alpha-ix)^{-1}$, the support principle and \eqref{HZ-1.5-res} imply
\[
  i(i\alpha I+T)^{-1}
  =\int_{\R}f_\alpha(t)U_t\,dt.
  \tag{HZ-op-ext}\label{HZ-op-ext}
\]
Therefore
\[
\begin{aligned}
  \norm{(i\alpha I+T)^{-1}}
  &\le \int_{\R}\abs{f_\alpha(t)}\norm{U_t}\,dt\\
  &=\norm{f_\alpha}_1
   =\frac1\alpha\tanh\!\left(\frac{\pi\alpha}{2}\right).
\end{aligned}
\tag{HZ-op-bound}\label{HZ-op-bound}
\]
Applying the same argument to $-T$ treats $\alpha<0$. Resolvent continuity at $\alpha=0$ gives
\[
  \norm{T^{-1}}
  \le\lim_{\alpha\downarrow0}
      \frac1\alpha\tanh\!\left(\frac{\pi\alpha}{2}\right)
  =\frac\pi2.
  \tag{HZ-alpha0}\label{HZ-alpha0}
\]

\subsection{The sharp antiperiodic translation model}

Define $X$ to be the space of measurable complex functions modulo almost-everywhere equality satisfying
\[
  f(s+\pi)=-f(s),
  \qquad
  \norm f_X:=\int_0^\pi\abs{f(s)}\,ds<\infty.
  \tag{HZ-anti}\label{HZ-anti}
\]
Translations
\[
  (U_tf)(s):=f(s-t)
  \tag{HZ-shift}\label{HZ-shift}
\]
form a strongly continuous isometric group. Let $iT$ be its generator. Since $U_\pi=-I$, splitting the Laplace integral into intervals of length $\pi$ gives
\[
  (\alpha I-iT)^{-1}
  =\int_0^\infty e^{-\alpha t}U_t\,dt
  =\frac{1}{1+e^{-\alpha\pi}}
     \int_0^\pi e^{-\alpha t}U_t\,dt,
  \qquad\alpha>0.
  \tag{HZ-periodic-res}\label{HZ-periodic-res}
\]
\sourcecomment{The printed display has $1+e^{-\alpha t}$ in the denominator. The next display, the period splitting, and the geometric-series derivation force $1+e^{-\alpha\pi}$; this note records the corrected identity.}

The right side is analytic in $\alpha\in\C\setminus i(2\Z+1)$. Thus $\sigma(T)\subset2\Z+1$. Conversely, for every odd integer $n$, the function $s\mapsto e^{-ins}$ belongs to $X$ and is an eigenvector of $T$ with eigenvalue $n$. Hence
\[
  \sigma(T)=2\Z+1.
  \tag{HZ-odd-spectrum}\label{HZ-odd-spectrum}
\]
The upper bound becomes
\[
  \norm{(\alpha I-iT)^{-1}}
  \le \frac{1}{1+e^{-\alpha\pi}}
      \int_0^\pi e^{-\alpha t}\,dt
  =\frac1\alpha\tanh\!\left(\frac{\pi\alpha}{2}\right).
  \tag{HZ-sharp-upper}\label{HZ-sharp-upper}
\]

For $n\ge1$, define $f_n$ on $(0,\pi)$ by
\[
  f_n(s)=n\mathbf1_{[0,1/n)}(s)
  \tag{HZ-testfn}\label{HZ-testfn}
\]
and extend antiperiodically. Then $\norm{f_n}_X=1$. For $s\in[1/n,\pi)$,
\[
\begin{aligned}
  [(\alpha I-iT)^{-1}f_n](s)
  &=\frac{n}{1+e^{-\alpha\pi}}
      \int_{s-1/n}^{s}e^{-\alpha t}\,dt\\
  &=\frac{e^{-\alpha s}}{1+e^{-\alpha\pi}}
      \frac{n}{\alpha}(e^{\alpha/n}-1).
\end{aligned}
\tag{HZ-test-image}\label{HZ-test-image}
\]
Consequently,
\[
\begin{aligned}
 \norm{(\alpha I-iT)^{-1}}
 &\ge\limsup_{n\to\infty}
      \int_{1/n}^{\pi}
      \abs{[(\alpha I-iT)^{-1}f_n](s)}\,ds\\
 &=\frac1\alpha\tanh\!\left(\frac{\pi\alpha}{2}\right).
\end{aligned}
\tag{HZ-sharp-lower}\label{HZ-sharp-lower}
\]
Thus equality holds for $\alpha\ne0$, and continuity gives $\norm{T^{-1}}=\pi/2$.

\begin{corollary}[Source Corollary 1.6: Partington's conjecture]
For every Hermitian operator $T$ and every $\lambda\in\C\setminus\sigma(T)$,
\[
  \norm{(\lambda I-T)^{-1}}
  \le \frac\pi2\,\dist(\lambda,\sigma(T))^{-1}.
  \tag{HZ-1.6}\label{HZ-1.6}
\]
The constant $\pi/2$ is optimal.
\end{corollary}

\begin{proofoutline}
\item If $\dist(\lambda,\sigma(T))=\abs{\Im\lambda}$, the defining Hermitian resolvent estimate \eqref{HZ-H1} is already stronger.
\item Otherwise set
\[
  a:=\sqrt{\dist(\lambda,\sigma(T))^2-(\Im\lambda)^2}>0.
  \tag{HZ-normal-a}\label{HZ-normal-a}
\]
Translate by $\Re\lambda$ and scale by $a$. The normalized operator has spectrum outside $(-1,1)$ with at least one of $\pm1$ in its spectrum, and the normalized spectral parameter is $-i\alpha$ where
\[
  \alpha:=a^{-1}\Im\lambda.
  \tag{HZ-normal-alpha}\label{HZ-normal-alpha}
\]
\item Theorem 1.5 reduces the claim to
\[
  \frac1\alpha\tanh\!\left(\frac{\pi\alpha}{2}\right)
  \le \frac\pi2(1+\alpha^2)^{-1/2}
  \qquad(\alpha>0).
  \tag{HZ-hyperbolic-target}\label{HZ-hyperbolic-target}
\]
\item Differentiate both sides after normalizing:
\[
  \frac2\pi\tanh\!\left(\frac{\pi\alpha}{2}\right)
  =\int_0^\alpha\frac{du}{\cosh^2(\pi u/2)},
  \tag{HZ-hyperbolic-int}\label{HZ-hyperbolic-int}
\]
\[
  \frac{\alpha}{\sqrt{1+\alpha^2}}
  =\int_0^\alpha\frac{du}{(1+u^2)^{3/2}}.
  \tag{HZ-algebraic-int}\label{HZ-algebraic-int}
\]
The Taylor estimate
\[
  \cosh^2\!\left(\frac{\pi u}{2}\right)
  \ge\left(1+\frac{\pi^2u^2}{8}\right)^2
  \ge(1+u^2)^2
  \ge(1+u^2)^{3/2}
  \tag{HZ-denom-compare}\label{HZ-denom-compare}
\]
proves the integrand comparison and hence \eqref{HZ-hyperbolic-target}.
\end{proofoutline}

\section{Abstract minimal extrapolation on a locally compact abelian group}

\sourceanchor{Section 2.}

Let $G$ be a locally compact abelian group, $\widehat G$ its Pontryagin dual, and $M(G)$ the Banach space of finite complex regular Borel measures with total-variation norm $\norm\mu$.

Let $\varnothing\ne E\subset\widehat G$ be closed and let $\varphi:E\to\C$ be bounded and continuous. Assume at least one $\mu\in M(G)$ extrapolates $\varphi$:
\[
  \widehat\mu|_E=\varphi.
\]
The minimal extrapolation problem asks:
\begin{enumerate}[label=(\Roman*)]
\item compute
\[
  \inf\{\norm\mu:\mu\in M(G),\ \widehat\mu|_E=\varphi\};
  \tag{HZ-ME-I}\label{HZ-ME-I}
\]
\item decide whether the infimum is attained;
\item classify every minimizer.
\end{enumerate}

\begin{theorem}[Beurling, source Theorem 2.1]
Assume every nonempty relatively open subset of $E$ has nonzero Haar measure. Then
\[
  \{\mu\in M(G):\widehat\mu|_E=\varphi\}
  \tag{HZ-feasible}\label{HZ-feasible}
\]
is weak-star closed in $M(G)=C_0(G)^*$. If it is nonempty, the set of norm-minimizing extrapolants is a nonempty weak-star compact convex subset of $M(G)$.
\end{theorem}

\subsection{Compact spectral complement forces absolute continuity}

Suppose $\widehat G\setminus E$ is relatively compact and $\mu,\nu\in M(G)$ agree on $E$. Then
\[
  \supp(\widehat\mu-\widehat\nu)
  \subset\widehat G\setminus E
  \tag{HZ-compact-support}\label{HZ-compact-support}
\]
is compact. Hence $\widehat\mu-\widehat\nu\in L^1(\widehat G)$, and Fourier inversion implies
\[
  \mu-\nu\in L^1(G)\cap C_0(G).
  \tag{HZ-ac-difference}\label{HZ-ac-difference}
\]
This is the step that converts measure minimizers into actual $L^1$ functions in the application $E=\R\setminus(-1,1)$.

\begin{corollary}[Source Corollary 2.2]
Assume $E$ satisfies the hypotheses of Theorem 2.1, $\widehat G\setminus E$ is relatively compact, and
\[
  \varphi=\widehat g|_E
  \qquad\text{for some }g\in L^1(G).
\]
Then every feasible measure belongs to
\[
  L^1(G)\cap(g+C_0(G)).
  \tag{HZ-feasible-L1}\label{HZ-feasible-L1}
\]
The minimizing set is a nonempty convex subset of this affine space and is compact for the weak topology induced by the pairing $L^1(G)\times C_0(G)$.
\end{corollary}

\begin{theorem}[Beurling--Sz\H{o}kefalvi-Nagy, source Theorem 2.3]
Let $\varnothing\ne E\subset\widehat G$ and $f\in L^1(G)$. The following are equivalent.
\begin{enumerate}[label=(\arabic*)]
\item $f$ is minimal among measure extrapolants:
\[
  \norm f_1
  =\inf\{\norm\mu:\mu\in M(G),\ \widehat\mu|_E=\widehat f|_E\}.
  \tag{HZ-2.3a}\label{HZ-2.3a}
\]
\item $f$ is minimal among $L^1$ extrapolants:
\[
  \norm f_1
  =\inf\{\norm g_1:g\in L^1(G),\ \widehat g|_E=\widehat f|_E\}.
  \tag{HZ-2.3b}\label{HZ-2.3b}
\]
\item There exists $\sigma\in L^\infty(G)$ such that
\[
  \norm\sigma_\infty\le1,
  \qquad \sigma f=\abs f\quad\text{a.e.},
  \tag{HZ-dual-phase}\label{HZ-dual-phase}
\]
and
\[
  \int_G\sigma(t)h(t)\,dt=0
  \quad\text{whenever }h\in L^1(G),\ \widehat h|_E=0.
  \tag{HZ-annihilator}\label{HZ-annihilator}
\]
\end{enumerate}
\end{theorem}

The function $\sigma$ is a norming functional for $f$ that annihilates the closed subspace of admissible perturbations. It is the exact dual certificate for minimality.

\begin{corollary}[Source Corollary 2.4]
Assume $E$ is closed, $\widehat G\setminus E$ is relatively compact, and $f\in L^1(G)$ admits a certificate $\sigma$ satisfying \eqref{HZ-dual-phase} and \eqref{HZ-annihilator}. Then:
\begin{enumerate}[label=(\arabic*)]
\item $f$ is a minimal measure extrapolant:
\[
  \norm f_1
  =\inf\{\norm\mu:\mu\in M(G),\ \widehat\mu|_E=\widehat f|_E\};
  \tag{HZ-2.4a}\label{HZ-2.4a}
\]
\item every feasible measure is $g(t)\,dt$ for some $g\in L^1(G)$, and every such $g$ satisfies
\[
  \norm f_1=\int_G\sigma(t)g(t)\,dt;
  \tag{HZ-2.4b}\label{HZ-2.4b}
\]
\item for feasible $g\in L^1(G)$,
\[
  \norm g_1=\norm f_1
  \quad\Longleftrightarrow\quad
  \sigma g=\abs g\quad\text{a.e.}
  \tag{HZ-2.4c}\label{HZ-2.4c}
\]
\end{enumerate}
\end{corollary}

\sourcecomment{The printed item (1) of Corollary 2.4 places $\mu$ in $M(\widehat G)$. The feasible measures throughout the section are measures on $G$, and the proof invokes Theorem 2.3 with $M(G)$; this note uses the mathematically consistent $M(G)$.}

\begin{argument}
\item Theorem 2.3 gives \eqref{HZ-2.4a}.
\item Corollary 2.2 gives $\mu=g(t)\,dt$. Since $\widehat{f-g}|_E=0$, annihilation yields
\[
  0=\int_G\sigma(f-g),
\]
and $\sigma f=\abs f$ gives \eqref{HZ-2.4b}.
\item If $\sigma g=\abs g$, then \eqref{HZ-2.4b} gives equality of norms. Conversely, if the norms agree, then
\[
  \int_G\bigl(\abs g-\Re(\sigma g)\bigr)=0.
\]
The integrand is nonnegative, so it vanishes a.e. Together with $\abs{\sigma g}\le\abs g$, this forces $\sigma g=\abs g$ a.e.
\end{argument}

\section{Minimal extrapolation on the real line}

\sourceanchor{Section 3. Here $G=\widehat G=\R$ and $E=\R\setminus(-1,1)$.}

\subsection{A periodic annihilator and the sign certificate}

Let $\sigma\in L^\infty(\R)$ be $2\pi$-periodic. If $h\in L^1(\R)$ and
\[
  \widehat h(k)=0\qquad(k\in\Z\setminus\{0\}),
  \tag{HZ-integer-zero}\label{HZ-integer-zero}
\]
then
\[
  \int_{\R}\sigma(t)h(t)\,dt
  =\frac1{2\pi}\left(\int_{-\pi}^{\pi}\sigma(t)\,dt\right)
       \left(\int_{\R}h(t)\,dt\right).
  \tag{HZ-periodic-pairing}\label{HZ-periodic-pairing}
\]
Indeed, it suffices by weak-star density of trigonometric polynomials to verify the identity for $\sigma(t)=e^{ikt}$, where it is exactly the Fourier coefficient condition.

If $\sigma$ has mean zero over one period, then it annihilates every $h$ with
\[
  \widehat h|_{\R\setminus(-1,1)}=0,
  \tag{HZ-bandlimited}\label{HZ-bandlimited}
\]
because all nonzero integers belong to the exterior set. The source chooses
\[
  \sigma(t)=\sgn(\sin t),
  \qquad\abs{\sigma(t)}=1\text{ a.e.},
  \qquad\int_{-\pi}^{\pi}\sigma(t)\,dt=0.
  \tag{HZ-sign-cert}\label{HZ-sign-cert}
\]

\begin{proposition}[Source Proposition 3.1]
Let $f\in L^1(\R)$ satisfy
\[
  f(t)\sin t\ge0\quad\text{a.e.}
  \tag{HZ-sign-condition}\label{HZ-sign-condition}
\]
Then:
\begin{enumerate}[label=(\arabic*)]
\item $f$ minimizes total variation among all measures with the same exterior Fourier data;
\item every feasible measure is $g(t)\,dt$ and
\[
  \norm f_1=\int_{\R}g(t)\sgn(\sin t)\,dt;
  \tag{HZ-3.1-pair}\label{HZ-3.1-pair}
\]
\item a feasible $g\in L^1$ is also minimizing exactly when
\[
  g(t)\sin t\ge0\quad\text{a.e.}
  \tag{HZ-3.1-min}\label{HZ-3.1-min}
\]
\end{enumerate}
\end{proposition}

This is Corollary 2.4 with the explicit dual certificate \eqref{HZ-sign-cert}.

\begin{proposition}[Source Proposition 3.2: uniqueness]
Let $f$ be continuous on $\R\setminus\{0\}$, integrable, and satisfy
\[
  f(t)\sin t\ge0\qquad(t\ne0).
\]
Then $f(t)\,dt$ is the unique minimum-norm measure whose Fourier transform agrees with $\widehat f$ on $\R\setminus(-1,1)$.
\end{proposition}

\begin{argument}
\item If another minimizer is $g(t)\,dt$, Proposition 3.1 gives $g(t)\sin t\ge0$ a.e.
\item Set $h=g-f$. Since $\supp\widehat h\subset[-1,1]$, Fourier inversion gives a continuous representative
\[
  h(t)=\frac1{2\pi}\int_{-1}^{1}\widehat h(x)e^{-ixt}\,dx.
  \tag{HZ-band-inversion}\label{HZ-band-inversion}
\]
Thus $g$ may be chosen continuous on $\R\setminus\{0\}$.
\item At every nonzero $n\in\Z$, the sign of $\sin t$ changes at $n\pi$. Continuity and the two sign conditions force
\[
  f(n\pi)=g(n\pi)=0.
  \tag{HZ-lattice-zero}\label{HZ-lattice-zero}
\]
Hence
\[
  \frac1{2\pi}\int_{-1}^{1}\widehat h(x)e^{-in\pi x}\,dx=0
  \qquad(n\in\Z\setminus\{0\}).
  \tag{HZ-Fourier-coeff-zero}\label{HZ-Fourier-coeff-zero}
\]
\item Uniqueness of Fourier coefficients in $L^1([-1,1])$ implies $\widehat h$ is a.e. constant on $[-1,1]$. Continuity of $\widehat h$ and its vanishing off $[-1,1]$ force that constant to be zero. Therefore $h=0$.
\end{argument}

\subsection{A Laplace-transform integrability lemma}

For $F:[0,\infty)\to\C$, write
\[
  \mathcal L_y(F)(t):=\int_0^\infty F(y)e^{-ty}\,dy
  \qquad(t>0).
  \tag{HZ-Laplace}\label{HZ-Laplace}
\]

\begin{lemma}[Source Lemma 3.3]
Let $F$ be bounded and continuous on $[0,\infty)$, with $F(0)=0$, and differentiable at $0$. Then
\[
  \int_0^\infty\abs{\sin t}\,
     \abs{\mathcal L_y(F)(t)}\,dt<\infty.
  \tag{HZ-3.3}\label{HZ-3.3}
\]
\end{lemma}

Since $F(y)=O(y)$ near $0$ and is bounded at infinity,
\[
  c:=\sup_{y>0}\frac{\abs{F(y)}}{1-e^{-y}}<\infty.
  \tag{HZ-cF}\label{HZ-cF}
\]
Therefore
\[
  \abs{\mathcal L_y(F)(t)}
  \le c\int_0^\infty(1-e^{-y})e^{-ty}\,dy
  =\frac{c}{t(1+t)},
  \tag{HZ-Laplace-bound}\label{HZ-Laplace-bound}
\]
and \eqref{HZ-3.3} follows from integrability of $\abs{\sin t}/(t(1+t))$.

\section{The explicit minimizer for nonzero parameter}

\sourceanchor{Lemma 3.4.}

Fix $\alpha>0$. Let
\[
  g_\alpha(t):=e^{-\alpha t}\mathbf1_{[0,\infty)}(t),
  \qquad
  \widehat g_\alpha(x)=\frac1{\alpha-ix}.
  \tag{HZ-galpha}\label{HZ-galpha}
\]

\begin{lemma}[Source Lemma 3.4]
There is exactly one $f_\alpha\in L^1(\R)$ satisfying
\[
  f_\alpha(t)\sin t\ge0\quad\text{a.e.}
  \tag{HZ-3.4-sign}\label{HZ-3.4-sign}
\]
and
\[
  \widehat f_\alpha(x)=\frac1{\alpha-ix}
  \qquad(|x|\ge1).
  \tag{HZ-3.4-data}\label{HZ-3.4-data}
\]
It is given, for $t\ne0$, by
\[
 f_\alpha(t)=
 \begin{cases}
 \displaystyle
 \sin t\int_0^\infty
 \left(\frac1{e^{\pi\alpha}+1}
       -\frac1{e^{\pi(\alpha+y)}+1}\right)e^{ty}\,dy,
 &t<0,\\[2.2ex]
 \displaystyle
 \sin t\int_0^\infty
 \left(\frac1{e^{\pi(\alpha-y)}+1}
       -\frac1{e^{\pi\alpha}+1}\right)e^{-ty}\,dy,
 &t>0.
 \end{cases}
 \tag{HZ-3.4-formula}\label{HZ-3.4-formula}
\]
Its Fourier transform on all of $\R$ is
\[
 \widehat f_\alpha(x)
 =\frac1{\alpha-ix}
 -\pi\left(
     \frac1{e^{\pi\alpha}+1}
     +\frac1{e^{\pi(\alpha-ix)}-1}
   \right)\mathbf1_{[-1,1]}(x).
 \tag{HZ-3.4-Fourier}\label{HZ-3.4-Fourier}
\]
\end{lemma}

\subsection{Uniqueness and reconstruction of the compactly supported correction}

Suppose $f_\alpha$ satisfies the two conditions and set
\[
  h_\alpha:=g_\alpha-f_\alpha.
\]
Then
\[
  \supp\widehat h_\alpha\subset[-1,1],
  \tag{HZ-ha-support}\label{HZ-ha-support}
\]
and $h_\alpha$ has a continuous representative. The sign condition forces
\[
  f_\alpha(n\pi)=0\qquad(n\in\Z\setminus\{0\}).
\]
Thus, for $n\ne0$,
\[
  \frac1{2\pi}\int_{-1}^{1}\widehat h_\alpha(x)e^{-in\pi x}\,dx
  =h_\alpha(n\pi)
  =
  \begin{cases}
    e^{-n\pi\alpha},&n>0,\\
    0,&n<0.
  \end{cases}
  \tag{HZ-3.4-coeff}\label{HZ-3.4-coeff}
\]
The resulting absolutely convergent Fourier series on $[-1,1]$ is
\[
  \widehat h_\alpha(x)
  =\pi\left(c_\alpha+\sum_{n=1}^\infty
       e^{-n\pi\alpha}e^{in\pi x}\right)
  =\pi\left(c_\alpha+
       \frac1{e^{\pi(\alpha-ix)}-1}\right).
  \tag{HZ-3.4-series}\label{HZ-3.4-series}
\]
Continuity and support imply $\widehat h_\alpha(\pm1)=0$, so
\[
  c_\alpha=\frac1{e^{\pi\alpha}+1}.
  \tag{HZ-calpha}\label{HZ-calpha}
\]
Therefore
\[
  h_\alpha(t)
  =\frac12\int_{-1}^{1}
   \left(
    \frac1{e^{\pi\alpha}+1}
    +\frac1{e^{\pi(\alpha-ix)}-1}
   \right)e^{ixt}\,dx.
  \tag{HZ-3.4-h}\label{HZ-3.4-h}
\]
This identifies the only possible candidate $f_\alpha=g_\alpha-h_\alpha$.

\subsection{Contour calculation for \texorpdfstring{$t<0$}{t less than 0}}

Introduce
\[
  F_t(z):=\frac12\left(
    \frac1{e^{\pi\alpha}+1}
    +\frac1{e^{\pi(\alpha-iz)}-1}
  \right)e^{-itz}.
  \tag{HZ-meromorphic}\label{HZ-meromorphic}
\]
Its poles are $z=-i\alpha+2n$, $n\in\Z$, with
\[
  \operatorname*{Res}_{z=-i\alpha}F_t(z)
  =-\frac{1}{2\pi i}e^{-\alpha t}.
  \tag{HZ-residue}\label{HZ-residue}
\]
For $t<0$, integrate around
\[
  R(T)=\{z:-1\le\Re z\le1,\ 0\le\Im z\le T\}.
\]
There are no poles, so
\[
  \int_{\partial R(T)}F_t(z)\,dz=0.
  \tag{HZ-contour-up}\label{HZ-contour-up}
\]
The top edge vanishes as $T\to\infty$, and the vertical edges yield
\[
  h_\alpha(t)
  =-\sin t\int_0^\infty
    \left(\frac1{e^{\pi\alpha}+1}
          -\frac1{e^{\pi(\alpha+y)}+1}
    \right)e^{ty}\,dy.
  \tag{HZ-v}\label{HZ-v}
\]
Since $g_\alpha(t)=0$ for $t<0$, this gives the first branch of \eqref{HZ-3.4-formula}.

\subsection{Contour calculation for \texorpdfstring{$t>0$}{t greater than 0}}

For $t>0$, integrate around
\[
  R'(T)=\{z:-1\le\Re z\le1,\ -T\le\Im z\le0\}.
\]
When $T>\alpha$, the contour contains exactly the pole $-i\alpha$, and the residue theorem gives
\[
  \int_{\partial R'(T)}F_t(z)\,dz=e^{-\alpha t}.
  \tag{HZ-contour-down}\label{HZ-contour-down}
\]
Again the horizontal edge vanishes, and the vertical edges give
\[
  h_\alpha(t)
  =e^{-\alpha t}
   +\sin t\int_0^\infty
    \left(\frac1{e^{\pi\alpha}+1}
          -\frac1{e^{\pi(\alpha-y)}+1}
    \right)e^{-ty}\,dy.
  \tag{HZ-vii-h}\label{HZ-vii-h}
\]
Subtracting from $g_\alpha(t)=e^{-\alpha t}$ gives the second branch of \eqref{HZ-3.4-formula}.

The integrands in both branches have the sign needed to show
\[
  f_\alpha(t)\sin t\ge0\qquad(t\ne0).
  \tag{HZ-3.4-sign-pointwise}\label{HZ-3.4-sign-pointwise}
\]
Lemma 3.3 proves $f_\alpha\in L^1$. Fourier inversion applied to \eqref{HZ-3.4-h} then yields \eqref{HZ-3.4-Fourier} and the exterior data \eqref{HZ-3.4-data}. Uniqueness also follows immediately from Proposition 3.2.

\section{The exact norm of the minimizer}

\sourceanchor{Lemma 3.5.}

\begin{lemma}[Source Lemma 3.5]
For $\alpha>0$,
\[
  \norm{f_\alpha}_1
  =\frac1\alpha\tanh\!\left(\frac{\pi\alpha}{2}\right).
  \tag{HZ-3.5}\label{HZ-3.5}
\]
\end{lemma}

For $0<t<\pi$, define the antiperiodization
\[
  k_\alpha(t):=\sum_{n\in\Z}(-1)^n f_\alpha(t+n\pi).
  \tag{HZ-kalpha}\label{HZ-kalpha}
\]
The sum converges a.e. and in $L^1(0,\pi)$. Because $f_\alpha(t)\sin t\ge0$,
\[
  \norm{f_\alpha}_1
  =\int_{\R}f_\alpha(t)\sgn(\sin t)\,dt
  =\int_0^\pi k_\alpha(t)\,dt.
  \tag{HZ-norm-antiperiod}\label{HZ-norm-antiperiod}
\]
For every odd integer $m$,
\[
  \int_0^\pi k_\alpha(t)e^{imt}\,dt
  =\int_{\R}f_\alpha(t)e^{imt}\,dt
  =\frac1{\alpha-im}.
  \tag{HZ-viii}\label{HZ-viii}
\]
Apply the same construction to $g_\alpha=e^{-\alpha t}\mathbf1_{[0,\infty)}$. Its antiperiodization is
\[
  \ell_\alpha(t)
  :=\sum_{n\in\Z}(-1)^n g_\alpha(t+n\pi)
  =\sum_{n=0}^\infty(-1)^n e^{-\alpha(t+n\pi)}
  =\frac{e^{-\alpha t}}{1+e^{-\alpha\pi}}.
  \tag{HZ-lalpha}\label{HZ-lalpha}
\]
The source prints the denominator in the last expression in an algebraically equivalent but OCR-prone form; the geometric series fixes it unambiguously.

The odd, $2\pi$-periodic extensions of $k_\alpha$ and $\ell_\alpha$ have the same Fourier coefficients, so they agree a.e. on $(0,\pi)$. Therefore
\[
\begin{aligned}
  \norm{f_\alpha}_1
  &=\int_0^\pi\frac{e^{-\alpha t}}{1+e^{-\alpha\pi}}\,dt\\
  &=\frac1\alpha\frac{1-e^{-\alpha\pi}}{1+e^{-\alpha\pi}}
   =\frac1\alpha\tanh\!\left(\frac{\pi\alpha}{2}\right).
\end{aligned}
\tag{HZ-norm-computation}\label{HZ-norm-computation}
\]

\section{The complete minimal-extrapolation theorem}

\sourceanchor{Theorem 3.6.}

\begin{theorem}[Source Theorem 3.6]
For every $\alpha\in\R$ there is a unique measure $\nu_\alpha\in M(\R)$ such that
\[
  \widehat\nu_\alpha(x)=\frac1{\alpha-ix}
  \qquad(|x|\ge1)
  \tag{HZ-3.6-data}\label{HZ-3.6-data}
\]
and
\[
  \norm{\nu_\alpha}
  =\inf\left\{\norm\mu:
      \mu\in M(\R),\ 
      \widehat\mu(x)=\frac1{\alpha-ix}\text{ for }|x|\ge1
    \right\}.
  \tag{HZ-3.6-min}\label{HZ-3.6-min}
\]
The minimizer is absolutely continuous:
\[
  d\nu_\alpha(t)=f_\alpha(t)\,dt.
\]
For $t\ne0$,
\[
 f_\alpha(t)=
 \begin{cases}
 \displaystyle
 \frac{\sin t}{2}\int_0^\infty
 \left(
  \tanh\!\frac{\pi(\alpha+y)}2-
  \tanh\!\frac{\pi\alpha}2
 \right)e^{ty}\,dy,
 &t<0,\\[2.2ex]
 \displaystyle
 \frac{\sin t}{2}\int_0^\infty
 \left(
  \tanh\!\frac{\pi\alpha}2-
  \tanh\!\frac{\pi(\alpha-y)}2
 \right)e^{-ty}\,dy,
 &t>0.
 \end{cases}
 \tag{HZ-3.6-f}\label{HZ-3.6-f}
\]
For $x\ne0$,
\[
  \widehat f_\alpha(x)
  =\frac1{\alpha-ix}
   +\frac\pi2\left(
      \tanh\!\frac{\pi\alpha}{2}
      -\coth\!\frac{\pi(\alpha-ix)}{2}
    \right)\mathbf1_{[-1,1]}(x).
  \tag{HZ-3.6-Fourier}\label{HZ-3.6-Fourier}
\]
Finally,
\[
  \norm{\nu_\alpha}=\norm{f_\alpha}_1
  =
  \begin{cases}
    \displaystyle\frac1\alpha\tanh\!\left(\frac{\pi\alpha}{2}\right),
      &\alpha\ne0,\\[1.1ex]
    \displaystyle\frac\pi2,&\alpha=0.
  \end{cases}
  \tag{HZ-3.6-norm}\label{HZ-3.6-norm}
\]
\end{theorem}

\subsection{Negative parameters}

For $h:\R\to\C$, write $\check h(x)=h(-x)$. Since
\[
  \frac1{-\alpha-ix}
  =-\left(\frac1{\alpha-ix}\right)\!\check{\ },
  \tag{HZ-reflect-multiplier}\label{HZ-reflect-multiplier}
\]
uniqueness for $\alpha>0$ implies
\[
  f_{-\alpha}=-\check f_\alpha.
  \tag{HZ-reflect-density}\label{HZ-reflect-density}
\]
The oddness of $\tanh$ extends \eqref{HZ-3.6-f}, \eqref{HZ-3.6-Fourier}, and \eqref{HZ-3.6-norm} to $\alpha<0$.

\subsection{The limiting parameter \texorpdfstring{$\alpha=0$}{alpha equals 0}}

The source establishes pointwise continuity of $\alpha\mapsto f_\alpha(t)$ for $t\ne0$ and then proves a uniform integrable majorant. First,
\[
  \abs{f_\alpha(t)}\le\abs{\frac{\sin t}{t}}.
  \tag{HZ-xi}\label{HZ-xi}
\]
If $F$ and $F'$ are bounded with $F(0)=0$, integration by parts gives
\[
  \mathcal L_y(F)(t)=\frac1t\mathcal L_y(F')(t).
  \tag{HZ-Laplace-parts}\label{HZ-Laplace-parts}
\]
Applying this to the derivatives of the hyperbolic-tangent expressions yields
\[
 f_\alpha(t)=
 \begin{cases}
 \displaystyle
 -\frac{\pi\sin t}{4t}
  \int_0^\infty
  \frac{e^{ty}}{\cosh^2(\pi(\alpha+y)/2)}\,dy,
 &t<0,\\[2.2ex]
 \displaystyle
 \frac{\pi\sin t}{4t}
  \int_0^\infty
  \frac{e^{-ty}}{\cosh^2(\pi(\alpha-y)/2)}\,dy,
 &t>0.
 \end{cases}
 \tag{HZ-partial-integration-form}\label{HZ-partial-integration-form}
\]
Since $\cosh^{-2}\le1$,
\[
  \abs{f_\alpha(t)}\le\frac{\pi\abs{\sin t}}{4t^2}.
  \tag{HZ-tail-bound}\label{HZ-tail-bound}
\]
Combining \eqref{HZ-xi} and \eqref{HZ-tail-bound} gives an integrable bound independent of $\alpha$:
\[
  \abs{f_\alpha(t)}
  \le \frac{2\abs{\sin t}}{\abs t(1+\abs t)}.
  \tag{HZ-dominated}\label{HZ-dominated}
\]
Dominated convergence therefore proves continuity
\[
  \R\ni\alpha\longmapsto f_\alpha\in L^1(\R).
  \tag{HZ-L1-cont}\label{HZ-L1-cont}
\]
At $\alpha=0$,
\[
  f_0(t)
  =\frac{\sin t}{2}
     \int_0^\infty\tanh\!\left(\frac{\pi y}{2}\right)
       e^{-\abs t y}\,dy,
  \qquad t\ne0.
  \tag{HZ-f0}\label{HZ-f0}
\]
It is continuous away from $0$ and satisfies $f_0(t)\sin t\ge0$. Proposition 3.2 gives uniqueness, while $L^1$ continuity gives
\[
  \widehat f_0(x)=\frac1{-ix}\qquad(|x|\ge1)
  \tag{HZ-f0-transform}\label{HZ-f0-transform}
\]
and
\[
  \norm{f_0}_1
  =\lim_{\alpha\to0}
    \frac1\alpha\tanh\!\left(\frac{\pi\alpha}{2}\right)
  =\frac\pi2.
  \tag{HZ-f0-norm}\label{HZ-f0-norm}
\]
The source notes that this $\alpha=0$ case was already obtained by Sz\H{o}kefalvi-Nagy and Strausz.

\section{The exact reciprocal multiplier needed by the Sylvester theorem}

This section is a modern extraction from Theorem 3.6, not a separately numbered theorem in the source.

\begin{corollary}[Sharp reciprocal Fourier extension]
Define
\[
  g(t):=-i f_0(t).
  \tag{HZ-g0}\label{HZ-g0}
\]
Then
\[
  \widehat g(x)=\frac1x\qquad(|x|\ge1),
  \qquad
  \norm g_1=\frac\pi2.
  \tag{HZ-reciprocal-unit}\label{HZ-reciprocal-unit}
\]
For $\delta>0$, define
\[
  g_\delta(t):=g(\delta t).
  \tag{HZ-scale-density}\label{HZ-scale-density}
\]
Then
\[
  \widehat g_\delta(x)=\frac1x
  \qquad(\abs x\ge\delta),
  \qquad
  \norm{g_\delta}_1=\frac\pi{2\delta}.
  \tag{HZ-reciprocal-scaled}\label{HZ-reciprocal-scaled}
\]
Moreover, $\pi/(2\delta)$ is the least possible $L^1$ norm among all Fourier extensions of $x\mapsto1/x$ from $\abs x\ge\delta$.
\end{corollary}

\begin{argument}
\item By \eqref{HZ-f0-transform}, $\widehat f_0(x)=1/(-ix)=i/x$, so multiplication by $-i$ gives \eqref{HZ-reciprocal-unit}.
\item Under the convention \eqref{HZ-F},
\[
  \widehat{g(\delta\,\cdot)}(x)
  =\frac1\delta\widehat g(x/\delta).
  \tag{HZ-scale-FT}\label{HZ-scale-FT}
\]
For $\abs x\ge\delta$, this is $(1/\delta)(\delta/x)=1/x$. The $L^1$ norm scales by $1/\delta$.
\item Minimality follows by the inverse scaling map from Theorem 3.6 at $\alpha=0$.
\end{argument}

\begin{corollary}[Finite-dimensional Sylvester integral formula]
Let $A:F\to F$ and $B:E\to E$ be self-adjoint operators on finite-dimensional complex Hilbert spaces, and suppose
\[
  \abs{a-b}\ge\delta>0
  \qquad(a\in\sigma(A),\ b\in\sigma(B)).
  \tag{HZ-spectral-separation}\label{HZ-spectral-separation}
\]
If
\[
  AX-XB=C,
  \tag{HZ-Sylvester}\label{HZ-Sylvester}
\]
then
\[
  X=\int_{\R}g_\delta(t)e^{itA}Ce^{-itB}\,dt.
  \tag{HZ-Sylvester-integral}\label{HZ-Sylvester-integral}
\]
Consequently every unitarily invariant norm $N$ satisfies
\[
  \delta N(X)\le\frac\pi2N(C).
  \tag{HZ-Sylvester-bound}\label{HZ-Sylvester-bound}
\]
\end{corollary}

\begin{argument}
\item In eigenbases of $A$ and $B$, equation \eqref{HZ-Sylvester} is
\[
  (a_i-b_j)X_{ij}=C_{ij}.
  \tag{HZ-Sylvester-coeff}\label{HZ-Sylvester-coeff}
\]
\item Separation and \eqref{HZ-reciprocal-scaled} imply
\[
  \frac1{a_i-b_j}
  =\int_{\R}g_\delta(t)e^{it(a_i-b_j)}\,dt.
  \tag{HZ-coeff-integral}\label{HZ-coeff-integral}
\]
Substitution proves \eqref{HZ-Sylvester-integral} coefficientwise.
\item Since $e^{itA}$ and $e^{-itB}$ are unitary,
\[
\begin{aligned}
  N(X)
  &\le\int_{\R}\abs{g_\delta(t)}
      N(e^{itA}Ce^{-itB})\,dt\\
  &=\norm{g_\delta}_1N(C)
   =\frac\pi{2\delta}N(C).
\end{aligned}
\tag{HZ-UI-bound}\label{HZ-UI-bound}
\]
\end{argument}

For the repository's barycentric formulation, write
\[
  g_\delta(t)=\abs{g_\delta(t)}\omega(t),
  \qquad\abs{\omega(t)}=1
  \quad\text{a.e. on }\{g_\delta\ne0\}.
\]
Absorb the phase $\omega(t)$ into the left unitary. Multiplying \eqref{HZ-Sylvester-integral} by $\delta$ produces a positive measure of total mass $\pi/2$ whose normalized integral lies in the closed real convex hull of the two-sided unitary orbit of $C$. This is precisely the analytic content expected by
\code{sylvester_barycentricOrbitRepresentation_of_spectralDistance}.
The remaining Lean-specific issue is the real-field descent: the paper and its Fourier phase argument are complex. For $\R$, one must either complexify and descend a conjugation-invariant average or replace scalar phases by real $2\times2$ rotations.

\section{Formalization map}

\begin{formalizationmap}
Hermitian definition & Banach-space generator/resolvent equivalence & Not presently a target; Mathlib's self-adjoint Hilbert operators are a narrower setting. \\
Theorems 1.1--1.2, Lemma 1.3, Corollary 1.4 & Historical progression of universal resolvent constants & Literature context only. \\
Theorem 1.5 & Sharp normalized resolvent bound and antiperiodic equality model & Not formalized; useful as an independent endpoint after Fourier infrastructure exists. \\
Corollary 1.6 & Universal sharp $\pi/2$ Hermitian resolvent estimate & Not formalized; broader Banach-space operator theory than the current finite Hilbert-space project. \\
Theorem 2.1 and Corollary 2.2 & Existence/compactness of minimum-norm extrapolants and automatic absolute continuity & Largely delegated to functional analysis and Fourier theory not yet assembled in this repository. \\
Theorem 2.3 and Corollary 2.4 & $L^1$ dual certificate and classification of minimizers & Missing reusable theorem layer; a plausible standalone \code{ForMathlib} target. \\
Propositions 3.1--3.2 & The $\sgn(\sin t)$ certificate and uniqueness from lattice zeros & Missing; these are the cleanest route to source-faithful minimality. \\
Lemma 3.3 & Weighted Laplace-transform integrability & Missing but elementary once the paper's Laplace notation is represented. \\
Lemma 3.4 & Explicit nonzero-parameter extremizer via Fourier series and contour integration & Missing; analytically heavy and not needed to close only the $\alpha=0$ Sylvester seam. \\
Lemma 3.5 & Exact norm by antiperiodization & Missing; source-faithful proof of the sharp constant for $\alpha>0$. \\
Theorem 3.6, $\alpha=0$ & Unique $L^1$ extension of $1/(-ix)$ with norm $\pi/2$ & Exact scalar analytic root needed by \code{DavisKahan/Alternative/FiniteDimensional/Sylvester/ContinuousLinearMapBridge.lean}. \\
Scaled reciprocal corollary & $\widehat g_\delta(x)=1/x$ off $(-\delta,\delta)$ and $\norm{g_\delta}_1=\pi/(2\delta)$ & Should be the first formal theorem extracted from this source. \\
Sylvester integral & $X=\int g_\delta(t)e^{itA}Ce^{-itB}\,dt$ & This is the intended proof of the remaining barycentric-orbit theorem over $\C$. \\
Real-field descent & Replace complex phases by a real convex-orbit certificate & Not in the source; remains an explicit formalization obligation. \\
\end{formalizationmap}

\section{Recommended formalization decomposition}

A source-faithful implementation should not begin with contour integration. The shortest path to the active theorem seam is:
\begin{enumerate}
\item State an abstract reciprocal-extension certificate:
\[
  \exists g\in L^1(\R),\quad
  \widehat g(x)=x^{-1}\ (\abs x\ge1),\quad
  \norm g_1\le\pi/2.
\]
\item Prove the scaling theorem \eqref{HZ-reciprocal-scaled}.
\item Prove the finite spectral-coordinate operator integral \eqref{HZ-Sylvester-integral}.
\item Convert the integral into closed convex-hull membership of the two-sided unitary orbit and then into the existing finite orbit certificate.
\item Handle $\R$ by complexification/descent or real rotations.
\item Only then reconstruct Theorem 3.6 source-faithfully from Propositions 3.1--3.2, Lemmas 3.3--3.5, and the limiting argument.
\end{enumerate}

The logical dependency should remain visible: Theorem 3.6 supplies the sharp scalar measure; operator spectral calculus transports it to a two-sided unitary average; convex geometry and Fan dominance transport that average to all unitarily invariant norms.

\section{Source-faithfulness and correction ledger}

\begin{itemize}
\item The source proof order is preserved: inherited resolvent estimates, normalized sharp theorem, abstract minimal extrapolation, real-line sign certificate, explicit extremizer, exact norm, then the complete theorem.
\item No later Bhatia--Davis--McIntosh or textbook proof is attributed to Haagerup--Zsid\'o. The Sylvester corollary is labeled as a modern extraction from their scalar theorem.
\item The paper works with complex Banach-space Hermitian operators. The current Lean application uses finite-dimensional real or complex Hilbert spaces; these are distinct scopes.
\item The denominator in the antiperiodic resolvent identity is corrected from the printed $1+e^{-\alpha t}$ to $1+e^{-\alpha\pi}$.
\item Corollary 2.4's measure space is corrected from the printed $M(\widehat G)$ to $M(G)$.
\item The source's Fourier convention is retained exactly. This matters for the factor $-i$ turning $1/(-ix)$ into $1/x$.
\item Operator-valued integrals are stated with the topology actually used by the paper. In the finite-dimensional Sylvester extraction they may be treated as Bochner integrals.
\item The source proves existence, uniqueness, and exact minimal norm, not merely an upper bound. The distilled statement preserves all three.
\item The $\alpha=0$ case is historically credited by the source to Sz\H{o}kefalvi-Nagy and Strausz; this note preserves that attribution.
\end{itemize}

\section{Bibliographic dependency ledger}
\begingroup\small

The paper explicitly depends on the following prior mathematical inputs:
\begin{itemize}
\item Lumer--Phillips and Stampfli--Williams for the Banach-space Hermitian/generator framework;
\item Partington for the original conjecture, the real-axis sharp bound, and subharmonic improvements;
\item Zsid\'o for the earlier general resolvent-growth theorem;
\item Beurling, Sz\H{o}kefalvi-Nagy, Domar, Esseen, and Shapiro for minimal Fourier extrapolation and its dual certificate;
\item Evans for the support of the integrated one-parameter-group functional calculus;
\item Rudin for Fourier inversion on locally compact abelian groups;
\item Sz\H{o}kefalvi-Nagy--Strausz for the $\alpha=0$ extremal problem.
\end{itemize}
These are dependencies of the source proof, not claims that the present repository has distilled each cited work.
\endgroup

\end{document}
